% TOP_PROBLEM: {"id": 2516, "title": "Kourovka Notebook Problem 21.7", "category_id": 17, "category": {"id": 17, "name": "group_theory", "display_name": "Group Theory", "description": "Problems about groups, group actions, representations, and related algebraic structures.", "slug": "group-theory", "order_index": 17, "created_at": "2026-07-31T15:26:25.671Z"}, "status": "open", "problem_number": "KOU-21.7", "set_id": 10}
% FIRST_POSTED: 2026-10-01
% ENTRY_KIND: statement_only
% UnsolvedMath ID 2516; source catalogue code KOU-21.7
% !TeX program = lualatex
% Definition and sources only; this file is not an LLM solution attempt.
\documentclass[11pt,a4paper]{article}
\usepackage[margin=25mm]{geometry}
\usepackage{fontspec}
\setmainfont{lmroman10-regular.otf}[BoldFont=lmroman10-bold.otf,
 ItalicFont=lmroman10-italic.otf,BoldItalicFont=lmroman10-bolditalic.otf]
\usepackage{amsmath,amssymb,mathtools}
\usepackage{unicode-math}
\setmathfont{latinmodern-math.otf}
\AtBeginDocument{\let\setminus\smallsetminus}
\usepackage{xurl,hyperref}
\hypersetup{colorlinks=true,urlcolor=blue}
\setcounter{secnumdepth}{0}
\setlength{\parindent}{0pt}
\setlength{\parskip}{5pt}
\setlength{\emergencystretch}{3em}
\begin{document}
\section*{Kourovka Notebook Problem 21.7}\label{2516}
{\small UnsolvedMath ID 2516 · KOU-21.7 · Group Theory · Kourovka Notebook - New Problems, Issue 21}
\subsection{Definitions and mathematical statement}
A finite left brace is a finite set $B$ with operations $+$ and $\circ$
such that $(B,+)$ is an abelian group, $(B,\circ)$ is a group, and
\[
 a\circ(b+c)=(a\circ b)-a+(a\circ c)
                       \qquad(a,b,c\in B),
\]
where $-a$ is the inverse in the additive group. The additive identity
is also the multiplicative identity. No right distributive law is
assumed. A finite group is called an IYB-group if it is isomorphic to
the multiplicative group $(B,\circ)$ of some finite left brace.
This brace definition fixes the equivalent formulation used in the source.

For a finite group $G$, a Sylow $p$-subgroup is a subgroup of order the
largest power of $p$ dividing $|G|$. The group $G$ is soluble if its
derived series, formed by repeatedly taking the subgroup generated by
commutators, terminates at $1$.

Suppose $G$ is a finite soluble group and every Sylow subgroup of $G$
is an IYB-group. Must $G$ be an IYB-group? Equivalently, can one put an
abelian addition on the underlying set of $G$ so that its given group
multiplication and that addition satisfy the brace identity?

The brace structures assumed to exist on the Sylow subgroups may be
different and need not agree on intersections or be compatible with
conjugation in $G$. The conclusion asks for one brace structure on the
whole group, not merely separate structures on its prime-power subgroups.
\subsection{Short English statement}
If every Sylow subgroup of a finite soluble group is the multiplicative group of a left brace, must the whole group admit such a structure?
\subsection{Sources}
\begin{itemize}
\item[S1] ULAM.AI and the UnsolvedMath Contributors, supplied problems.json, record 2516 (KOU-21.7), Kourovka Notebook - New Problems, Issue 21. Catalogue identity and source transcription; CC BY 4.0.
\url{https://huggingface.co/datasets/ulamai/UnsolvedMath}.
\item[S2] The Kourovka Notebook, 21st edition, new problems, Problem 21.7. Expanded formulation of the supplied catalogue statement.
\url{https://arxiv.org/pdf/1401.0300v47}.
\end{itemize}
\end{document}
